\documentclass{elsarticle}
\usepackage[utf8]{inputenc}
\usepackage[T1]{fontenc}
\usepackage{geometry}
\usepackage{graphicx}
\usepackage{float}
\usepackage{hyperref}
\usepackage[english,indonesian]{babel}

\title{Optimasi Pemrosesan Template LaTeX Menggunakan Pendekatan Berbasis AI}
\author{Academic Support Agent}
\date{\today}

\begin{document}
\maketitle

\begin{abstract}

\end{abstract}

\chapter{PENDAHULUAN}


\chapter{TINJAUAN PUSTAKA}


% Original content (commented out):
% \documentclass[times,5p, twocolumn]{elsarticle}
\usepackage[spanish, es-tabla]{babel}
%\usepackage[latin1]{inputenc}   % Lengua latina
\usepackage{inputenc}   % Lengua latin
\usepackage{amsmath,amsthm,amssymb}
%\usepackage{flushend}           % Para igualar las columnas de la última página
\usepackage{graphicx}
%\usepackage[pdftex]{graphicx}
\usepackage{epstopdf}                           % Convertir imagenes eps a pdf compilando con pdfLatex
\epstopdfsetup{outdir=./}
\usepackage[final]{pdfpages}                    % Paquete para incluir archivos PDF como si traslaparas hojas
\DeclareGraphicsExtensions{.eps}
\usepackage{listings}
\usepackage{color}
\usepackage[figuresright]{rotating}
\usepackage{indentfirst}
\usepackage{parskip}
\spanishdecimal{.}
\usepackage{pslatex}
\usepackage{listings, color}
\usepackage{color}
\usepackage{xcolor}
\definecolor{dkgreen}{rgb}{0,0.6,0}
\definecolor{dred}{rgb}{0.545,0,0}
\definecolor{dblue}{rgb}{0,0,0.545}
\definecolor{lgrey}{rgb}{0.9,0.9,0.9}
\definecolor{gray}{rgb}{0.4,0.4,0.4}
\definecolor{darkblue}{rgb}{0.0,0.0,0.6}
\lstdefinelanguage{ListCpp}{
	basicstyle=\footnotesize \ttfamily \color{black} \bfseries,   
	breakatwhitespace=false,       
	breaklines=true,               
	captionpos=t,                   
	commentstyle=\color{dkgreen},   
	deletekeywords={...},          
	escapeinside={\%*}{*)},                  
	frame=single,                  
	language=C++,                
	keywordstyle=\color{orange},  
	morekeywords={analogRead,analogWrite,getSpeed,getADC,MCP23S17,
		MCP3208,AD8804,SpiMotorsV4,SPI,Expander32,MCP,Slave,setQEI,
		begin,pinMode,digitalWrite,digitalRead,Serial,loop,setup,print,println}, 
	identifierstyle=\color{black},
	stringstyle=\color{blue},      
	numbers=left,                 
	numbersep=5pt,                  
	numberstyle=\tiny\color{black}, 
	rulecolor=\color{black},        
	showspaces=false,               
	showstringspaces=false,        
	showtabs=false,                
	stepnumber=1,                   
	tabsize=5,                     
	title=\lstname,                 
}

\begin{document}
\begin{frontmatter}
\title{Guía de Reportes de Prácticas\\ El Título No Debe Contener Más de 16 Palabras}

\author{Primer Autor}
\ead{primero@correo.correo.correo}

% Si necesita especificar una URL para el autor, consulte 'ejemplo latex RIAI.tex'
% \ead[url]{www.autor1.es}

\author{Segundo Autor}
\ead{segundo@correo.correo.correo}

\author{Tercero Autor}
\ead{tercero@correo.correo.correo}

\author{Asesor Autor}
\ead{asesor@correo.correo.correo}

\cortext[cor1]{Asesor, Profesor Investigador}

\address{División de Ingeniería Mecatrónica\\Instituto Tecnológico Superior de Huichapan\\Dom. Conocido s/n., El Saucillo, Municipio de Huichapan, Hidalgo, México}

\begin{abstract}
Todo reporte de práctica tiene como propósito dejar evidencia escrita del trabajo experimental que se realiza en los laboratorios, para cumplir con este fin un reporte debe contener puntos clave que permitan al lector duplicar y/o modificar los  resultados. En la presente guía para la redacción y elaboración de reporte de práctica se muestra lo mínimo necesario para llevar acabo esta tarea de manera correcta.

All practice's reports has the purpose to leave written evidence of experimental work done in laboratories, to satisfy this purpose, a report should contain key points that allow the reader to duplicate and / or modify the results. This guide presents the correctly way to prepare a practice's report with the minimum required to complete this task. 

\end{abstract}
\begin{keyword}
Reporte \sep Investigación \sep Guía \sep Práctica.
\end{keyword}
\end{frontmatter}

%%%%%%%%%%%%%%%%%%%%%%  Practice's report body  %%%%%%%%%%%%%%%%%%%%%%
\section{Introducción}
\noindent
Un reporte de práctica como con reporte de investigación se deben redactar con lenguaje técnico, es decir; el texto deberá escribirse en forma cuantitativa, evitar hacer uso de la palabras cualitativas como: más, menos, poco, mucho, entre otras, también se debe escribir en voz activa y no conjugar los verbos en infinitivo o progresivo.

Estos reportes deben contener los siguientes puntos:
\begin{itemize}
 \item Título de la práctica
 \item Resumen
 \item Abstract
 \item Introducción
 \item Objetivos
 \item Planteamiento del problema
 \item Marco teórico
 \item Materiales y métodos
 \item Desarrollo
 \item Resultados y análisis
 \item Conclusiones
 \item Referencias
\end{itemize} 

\section{Resumen y abstract}
El resumen se redacta al concluir el reporte de práctica puesto es que éste debe contener de manera breve los objetivos, el problema que se plantea, los resultados y las conclusiones que se obtuvieron al finalizar la práctica. El abstract es una traducción del resumen al Inglés.

Las \textit{keywords} son las palabras clave que definen la práctica, deben estar en orden alfabetico, se recomienda no usar las misma palabras que en el título.

\section{Introducción}
En esta sección se introduce al lector en el tema, se dan ejemplos de aplicaciones similares en otros lugares del mundo o en la industria. La introducción es un apartado breve en el que se describe de una manera general un panorama del trabajo que se va a presentar. Debe presentarse de manera resumida el alcance del trabajo y un resumen del mismo.

\section{Objetivos}
Los objetivos son una parte fundamental del reporte, este apartado se deben responder las preguntas: A) ¿Qué se va a hacer en la práctica?, B) ¿Cómo se realizará la práctica? y C) ¿Para qué realizará la práctica?. Todos los objetivos deben ser claros, empezar con un verbo en infinitivo y con no más de dos verbos.

Ejemplo: Verificar los cálculos matemáticos que se obtuvieron en clase al implementar una ley de control tipo PID de posición en un motor de corriente directa mediante el uso de un microcontrolador.

\section{Planteamiento del problema}
En esta sección se redacta de forma precisa la problemática que se encuentra en la práctica se acostumbra terminar con una pregunta que engloba al problema.


\section{Marco teórico}
En el marco teórico se presentan las herramientas teóricas que sustentan los principios que se van a utilizar en el desarrollo de la práctica. En esta sección es común presentar una revisión de literatura respecto al estado actual del caso de estudio, así como los principios matemáticos. 

Cuando se reporta una investigación, este apartado se convierte en el cuerpo del trabajo y contiene la información que se reúne.

\section{Materiales y métodos}

El la sección de marco teórico se presentaron las herramientas teóricas que sustentan el desarrollo y los resultados de la práctica, en esta sección se mencionan y describe el equipo y los materiales que se usan en la práctica. Se acostumbra hacer mención a las características técnicas.

\section{Desarrollo}

En esta sección se debe presentar de manera ordenada los pasos que se siguieron durante la experimentación. Esta sección se fortalece mediante el uso de imágenes y una descripción secuencial de los pasos a seguir. Es importante la redacción de esta sección, pues de ella depende la reproducibilidad del experimento que se realizó.

\section{Resultados y análisis}

En esta sección se presentan los resultados obtenidos en la experimentación. Es fuertemente reforzada mediante el uso de gráficas obtenidas de los resultados y la presentación de datos en tablas. El análisis se refiere a la comparación entre los resultados que se obtuvieron y los resultados que se esperan, también puede reforzarse esta sección con el uso de gráficas comparativas y datos estadísticos comunes (desviación estándar, error cuadrático medio, promedio, media, mediana, moda, entre otras)., sin embargo no siempre es necesario mostrar estos datos.

\section{Conclusiones}

\cite{Yeadon} Las conclusiones que se obtienen se deben relacionar directamente con los resultados del análisis que previamente se realizó. Las conclusiones deben ser respecto a los objetivos que se plantearon en la introducción y deben plasmar el aprendizaje que obtuvo por el alumno de la práctica que desarrolló. Se deben evitar la redacción redundante, por lo que se recomienda utilizar párrafos independientes para cada idea que expresa una conclusión particular e independiente.


\section{Referencias}
En esta sección se muestran las referencias que se consultaron para el desarrollo de la introducción y del marco teórico, para verificar que la información que se utilizó tiene respaldo científico. Existen diversas maneras para llevar a cabo la presentación de las referencias, se deja esta opción a criterio de quien realice el trabajo, siempre y cuando siga un estándar fijo para todas las referencias, se recomienda haga uso de diferentes fuentes de información.

\section{Elementos auxiliares del reporte}
\input{./imagenes.tex}
\input{./tablas.tex}
\input{./ecuaciones.tex}

\subsection{Código}
Para ingresar código al reporte solo es necesario copiar y pegar el código dentro de un lstlisting con el lenguaje Cpp:

\begin{lstlisting}[language=ListCpp,
caption={Código \LaTeX para insertar código en C/C++},
label = {codigo_1}]

aqui va el codigo

\end{lstlisting}


Ejemplo del código 1:

\begin{lstlisting}[language=ListCpp,
caption={Firmware para Arduino de la prueba de integración de módulos externos.},
label = {codigo1}]
#include<SPI.h>
#include<Expander32.h>
#include<SpiMotorsV4.h>
#include<MCP23S17.h>
#include<MCP3208.h>
#include<AD8804.h>

SpiMotorsV4 motor(11);
MCP in(0,9);
MCP out(0,32);
MCP3208 adc(15);
AD8804 dac(26);

void setup(){
	Slave.begin(40,41);
	motor.setQEI(MODE_X4_RESET_ON_MATCH,VELOCITY_PULSE_DIV_1,DIV_BY_1,4800);
	out.pinMode(5,LOW);
	out.pinMode(6,LOW);
	out.pinMode(7,LOW);
	out.pinMode(8,LOW);
	out.digitalWrite(5,LOW);
	out.digitalWrite(7,LOW);
	in.pinMode(13,HIGH);
	in.pinMode(16,HIGH);
	Serial.begin(9600);
}

boolean state=false;
float speed = 0;
float temp = 0;
void loop(){
	state=!state;
	digitalWrite(13,state);
	out.digitalWrite(8,state);
	speed = motor.getSpeed();
	temp = adc.analogRead(CH0)* 0.122100122;
	if (in.digitalRead(13)){
		dac.analogWrite(A1,230);  //B
		dac.analogWrite(A2,120);  //G
		dac.analogWrite(A3,150);  //R
		out.digitalWrite(6,HIGH);
		Serial.print("temp= ");
		Serial.print(temp);
		Serial.print(" C\t");
		Serial.print("speed= ");
		Serial.print(speed);
		Serial.println("\t");
	}else if (in.digitalRead(16)){
		out.digitalWrite(6,LOW);
		dac.analogWrite(A1,0);
		dac.analogWrite(A2,0);
		dac.analogWrite(A3,255);
	}else {
		out.digitalWrite(6,LOW);
		dac.analogWrite(A1,160);
		dac.analogWrite(A2,120);
		dac.analogWrite(A3,255);
	}
}
\end{lstlisting}
\bibliographystyle{plain}
\bibliography{reporte}	
\end{document} 

\end{document}